\chapter{Pelvic Pain}

\section*{Definition}
Chronic or acute pain localized to the lower abdominal region below the umbilicus, involving reproductive, urinary, gastrointestinal, or musculoskeletal structures.

\section*{When to See a Doctor}
See your doctor if:
\begin{itemize}
  \item Sudden severe pelvic pain, fever, vaginal bleeding in pregnancy, or pain with dizziness/syncope
\end{itemize}

\section*{How It Develops}
Pelvic pain in women most commonly arises from gynecological sources (endometriosis, adenomyosis, pelvic inflammatory disease, ovarian cysts) but can originate from gastrointestinal (IBS, diverticulitis), urological (interstitial cystitis), or musculoskeletal (pelvic floor dysfunction) systems. Visceral pain is poorly localized and often referred.

\section*{Causes}
\begin{itemize}
  \item Endometriosis
  \item Pelvic inflammatory disease (PID)
  \item Ovarian cysts or torsion
  \item Uterine fibroids (leiomyomas)
  \item Adenomyosis
  \item Irritable bowel syndrome (IBS)
  \item Interstitial cystitis (painful bladder syndrome)
  \item Pelvic floor dysfunction
  \item Ectopic pregnancy (emergent)
  \item Chronic prostatitis (in men)
\end{itemize}

\section*{Related Symptoms}
\begin{itemize}
  \item Lower back pain
  \item Dysmenorrhea (painful periods)
  \item Dyspareunia (painful sex)
  \item Urinary frequency
  \item Constipation
\end{itemize}

\section*{Medical Evaluation}
\begin{itemize}
  \item Comprehensive history: pain character (cyclic versus constant), relationship to menses, intercourse, urination, and bowel movements; obstetric and sexual history.
  \item Bimanual pelvic examination assessing cervical motion tenderness, adnexal masses, uterine enlargement, and uterosacral nodularity suggestive of endometriosis.
  \item Urine pregnancy test (mandatory in reproductive-age women), CBC, CRP, chlamydia and gonorrhoea NAAT, urinalysis with culture, and tumour markers (CA-125) if clinically indicated.
  \item Transvaginal pelvic ultrasound to evaluate for ovarian cysts, endometriomas, fibroids, adenomyosis, and hydrosalpinx; MRI or laparoscopy if ultrasound is inconclusive or endometriosis is strongly suspected.
\end{itemize}

\section*{Treatment Options}
\begin{itemize}
  \item Pelvic inflammatory disease, endometriosis, pregnancy-related causes, urinary disease, and gastrointestinal causes require diagnosis-specific treatment by a clinician.
  \item Treatment may include analgesia, hormonal therapy, antibiotics, pelvic-floor therapy, or specialist procedures depending on the confirmed cause.
  \item Referral to gynaecology for diagnostic laparoscopy, to gastroenterology if IBS is suspected, or to pelvic floor physiotherapy for myofascial pain.
\end{itemize}

\section*{Self-Care and Home Management}
\begin{itemize}
  \item Apply a heating pad or take a warm bath for 15--20 minutes to relax pelvic floor muscles and relieve cramping.
  \item An appropriate over-the-counter pain reliever may help if it is safe for you and used exactly as labelled; ask a clinician or pharmacist if you are pregnant or have kidney disease, ulcers, bleeding risk, or other medical conditions.
  \item Avoid constipation by increasing dietary fibre (25--30 g/day) and water intake; straining exacerbates pelvic pain.
  \item Practise mindfulness-based stress reduction and gentle stretching (yoga, pelvic tilts) to help manage chronic pelvic pain.
\end{itemize}

\section*{References}
\begin{thebibliography}{99}
\bibitem{pelvic_pain:pelvicpain1} Speer LM, Mushkbar S, Erbele T. ``Chronic pelvic pain in women.'' Am Fam Physician. 2016;93(5):380-387.
\bibitem{pelvic_pain:pelvicpain2} Dysfunctional uterine bleeding and chronic pelvic pain. In: Williams Gynecology. 4th ed. McGraw-Hill; 2020.
\bibitem{pelvic_pain:pelvicpain3} Royal College of Obstetricians and Gynaecologists. ``Chronic pelvic pain: initial management.'' Green-top Guideline No. 41. RCOG; 2012.
\bibitem{pelvic_pain:pelvicpain4} Howard FM. ``Chronic pelvic pain.'' Obstet Gynecol. 2003;101(3):594-611.
\bibitem{pelvic_pain:pelvicpain5} Steege JF, Siedhoff MT. ``Chronic pelvic pain.'' Obstet Gynecol. 2014;124(3):616-629.
\bibitem{pelvic_pain:pelvicpain6} Lamvu G, Carrillo J, Ouyang C, et al. ``Chronic pelvic pain in women: a review.'' JAMA. 2021;325(23):2381-2391.
\end{thebibliography}
